\SourceLeaf{388}%
\noindent mittelst paralleler Normalen auf die Kugel vom Radius 1 abgebildet und entspricht ihm dort ein Flächenelement der Grösse
\[
A(B'C''-C'B'')+B(C'A''-A'C'')+C(A'B''-B'A'')\,dt\,du=\Delta\,dt\,du,
\]
so ist das Krümmungsmass der Fläche im Punkte $t,u$:
\[
K=\frac{\Delta}{mm}.
\]
Vermöge der Gleichung 24 wird alsdann:
\begin{equation}\tag*{[25]}
K=-\frac1{mm}\left(\frac{\partial\partial\log m}{\partial t^2}+\frac{\partial\partial\log m}{\partial u^2}\right),
\end{equation}
$K$ lässt sich also allein durch $m$ und die Ableitungen davon ausdrücken, oder \textls[120]{das Krümmungsmass behält denselben Werth bei allen Umformungen der Fläche, die deren Linienelement}
\[
\BareRoot{\{mm(dt^2+du^2)\}}
\]
\textls[120]{unverändert lassen.}]
\Needspace{10\baselineskip}\begin{center}[3.]\end{center}
\noindent\hspace*{1.2em} Die Differentiation von $na$ und $n\alpha$ gibt:
\[
\begin{aligned}
na''+an''&=\frac{\partial\partial A}{\partial t\partial u}+C\frac{\partial\partial B}{\partial u^2}-B\frac{\partial\partial C}{\partial u^2}=A^{(3)}+CB^{(4)}-BC^{(4)}\\
na''+\alpha n'&=\frac{\partial\partial A}{\partial t\partial u}+B\frac{\partial\partial C}{\partial t^2}-C\frac{\partial\partial B}{\partial t^2}=A^{(3)}+BC^{(2)}-CB^{(2)}.
\end{aligned}
\]
[Mithin ist]
\[
an''-\alpha n'=C(B^{(2)}+B^{(4)})-B(C^{(2)}+C^{(4)}),
\]
[und ebenso:
\[
\begin{aligned}
bn''-\beta n'&=A(C^{(2)}+C^{(4)})-C(A^{(2)}+A^{(4)})\\
cn''-\gamma n'&=B(A^{(2)}+A^{(4)})-A(B^{(2)}+B^{(4)}).]
\end{aligned}
\]
Hieraus [folgt mit Hülfe von 7, 8 und 9]
\[
\begin{aligned}
mmn''&=\alpha(A^{(2)}+A^{(4)})+\beta(B^{(2)}+B^{(4)})+\gamma(C^{(2)}+C^{(4)})\\
mmn'&=a(A^{(2)}+A^{(4)})+b(B^{(2)}+B^{(4)})+c(C^{(2)}+C^{(4)}).
\end{aligned}
\]
Nun ist ferner:
\[
mmn'+mnm'=aA^{(2)}+bB^{(2)}+cC^{(2)}+\alpha A^{(3)}+\beta B^{(3)}+\gamma C^{(3)}.
\]
